% Generated draft: requires mathematical and source review.
\documentclass[11pt]{article}
\usepackage[margin=25mm]{geometry}
\usepackage{fontspec}
\setmainfont{Latin Modern Roman}
\usepackage{amsmath,amssymb,mathtools,unicode-math}
\setmathfont{Latin Modern Math}
\usepackage{enumitem,needspace,xurl,hyperref}
\setlength{\parindent}{0pt}
\setlength{\parskip}{5pt}
\setlength{\emergencystretch}{3em}
\newcommand{\N}{\mathbb{N}}
\newcommand{\Z}{\mathbb{Z}}
\newcommand{\Q}{\mathbb{Q}}
\newcommand{\R}{\mathbb{R}}
\newcommand{\C}{\mathbb{C}}
\newcommand{\F}{\mathbb{F}}
\newcommand{\A}{\mathbb{A}}
\newcommand{\PP}{\mathbb{P}}
\newcommand{\E}{\mathbb{E}}
\newcommand{\eps}{\varepsilon}
\newcommand{\dd}{\,\mathrm{d}}
\begin{document}
\section*{UnsolvedMath ID: 11000220}
\textbf{Source problem code:} catalogue entry
\subsection{Definitions and mathematical statement}
For a zero-order partition $\alpha'=(\alpha'_1,\ldots,\alpha'_n)$, the abelian-differential stratum $\mathcal H(\alpha')$ consists of translation surfaces $(X,\omega)$ with the specified zeros of $\omega$, with $\sum_i\alpha'_i=2g-2$. In local period coordinates, periods lie in $H^1(X,Z;\C)$, where $Z$ is the zero set. The kernel, or absolute-period, foliation fixes the image in $H^1(X;\C)$ while relative periods vary. Its local relative directions have complex dimension $n-1$ in the usual labeled-zero setting.

\textbf{Problem 13 (kernel foliation).} For the space $N$ constructed in the omitted preceding discussion, ask:
\begin{enumerate}
\item Is $N$ a complex-analytic orbifold? Is it a complex-algebraic orbifold?
\item Under what conditions does
\[
 \dim_{\C}N=\dim_{\C}O+n-m
\]
hold?
\item Under what conditions does $N$ coincide with the entire connected component of its enveloping stratum $\mathcal H(\alpha'_1,\ldots,\alpha'_n)$?
\end{enumerate}
The extracted record does not define $N$, $O$, or $m$, or state whether closures or kernel-leaf saturations are used. Although $n$ indexes the zeros of the target stratum, the role of $m$ cannot be certified. These missing constructions are indispensable: the displayed dimension formula is a proposed relation for the source's objects, not a general identity for arbitrary subsets or for every kernel-foliation leaf. In particular, the local relative dimension $n-1$ does not by itself justify the formula with $n-m$.

The source cites McMullen [McM3] for classification of $\mathrm{GL}_2(\R)$ orbit closures in $\mathcal H(1,1)$ and $\mathcal H(2)$, using homologous saddle connections. A saddle connection is a straight flat segment between zeros with no zero in its interior. It describes cutting a surface in $\mathcal H(1,1)$ along a pair joining the distinct zeros to obtain a two-torus description and use Ratner's theorem. This genus-two methodological context does not define the missing $N$ and $O$ or settle the three stated questions.
\subsection{Short English statement}
For the source's kernel-foliation construction, determine when its resulting space is analytic or algebraic, when the stated dimension formula holds, and when it fills a connected stratum component. The construction of the spaces and the parameter $m$ is missing.
\subsection{Sources}
\begin{itemize}
\item Local UnsolvedMath registry, permanent ID 11000220, source code catalogue entry; no proposer or proposal year supplied.
\item Farb, edited volume \emph{Problems on Mapping Class Groups and Related Topics}, Problem 13, PDF page 261; \url{https://www.math.uchicago.edu/~farb/papers/mcgbook.pdf}. The registry records PDF-font-structure extraction and access on 2026-07-29; the inventory distinguishes 329 labeled headings by source page.
\item McMullen [McM3] and Ratner's theorem are named in the source, without complete bibliography entries or exact theorem hypotheses in this record.
\end{itemize}
\end{document}
